% GENERATED FILE. Edit canonical lessons, then run npm run generate:books.
% canonical-source-hash: fnv1a64:2cc01682591bc176
% canonical-lessons: SA-C421-palli, SA-C421-upanagaram, SA-C421-catvaram, SA-C421-harmyam, SA-C421-sambhasate

\chapter{Hamlets and Suburbs}
\label{ch:sa-hamlets-and-suburbs}
\hlchaptermodality{421}

\begin{chapteropening}
\textbf{By the end of this chapter:} \emph{I can say where I live and chat with people there.}

Complete the last lesson of chapter 421: I can say where I live and chat with people there.
\end{chapteropening}

\section[\texorpdfstring{\sk{पल्ली} (pallī)}{पल्ली (pallī)}]{\sk{पल्ली} (pallī) — a hamlet}

\label{lesson:SA-C421-palli}



\begin{quote}
Before the new one: say the Sanskrit for comes upon, then the Sanskrit for refreshes.
\end{quote}

\subsection*{You'll want to know: \sk{पल्ली}}
\textbf{\sk{पल्ली}} — \emph{pallī} — \textquotedblleft{}a hamlet\textquotedblright{}.

\begin{quote}
\sk{पर्वतस्य} \sk{अधः} \sk{एका} \sk{पल्ली} \sk{अस्ति।}
\end{quote}

\begin{quote}
\emph{parvatasya adhaḥ ekā pallī asti.}
\end{quote}

\begin{quote}
Below the mountain there is a hamlet.
\end{quote}

A feminine noun.

\begin{cousinweb}
\textbf{\sk{पल्ली}} (\emph{pallī}) is the word MW gives for a small village; CDIAL, after T. Burrow, traces it to Dravidian, and DEDR 4018 sets Tamil \emph{paḷḷi} and Kannada \emph{haḷḷi}, \textquotedblleft{}a hamlet\textquotedblright{}, in one family, so it is usually taken as a loan from Dravidian (MW; CDIAL 7972; DEDR 4018).
\end{cousinweb}

\begin{grammarlens}[title={What you've built}]
One more word for getting around.
\end{grammarlens}

\subsection*{Your turn}
\begin{itemize}
\raggedright
  \item \emph{Say it:} \emph{pallī}
  \item \emph{Say it:} \emph{pallī}, once more
  \item \emph{Say it:} \emph{parvatasya adhaḥ ekā pallī asti.} — Below the mountain there is a hamlet.
  \item \emph{Recall:} say the Sanskrit for comes upon, then the Sanskrit for refreshes, then say \emph{pallī} again
\end{itemize}

\begin{tcolorbox}[breakable,colback=teal!4,colframe=teal!35!black,title={Before you move on}]
What does \sk{पल्ली} mean? (\textquotedblleft{}A hamlet\textquotedblright{}.) Say it once more.
\end{tcolorbox}
\section[\texorpdfstring{\sk{उपनगरम्} (upanagaram)}{उपनगरम् (upanagaram)}]{\sk{उपनगरम्} (upanagaram) — a suburb}

\label{lesson:SA-C421-upanagaram}



\begin{quote}
Before the new one: say the Sanskrit for refreshes, then the Sanskrit for a hamlet.
\end{quote}

\subsection*{You'll want to know: \sk{उपनगरम्}}
\textbf{\sk{उपनगरम्}} — \emph{upanagaram} — \textquotedblleft{}a suburb\textquotedblright{}.

\begin{quote}
\sk{मम} \sk{मित्रम्} \sk{उपनगरे} \sk{वसति।}
\end{quote}

\begin{quote}
\emph{mama mitram upanagare vasati.}
\end{quote}

\begin{quote}
My friend lives in a suburb.
\end{quote}

A neuter noun.

\begin{cousinweb}
\textbf{\sk{उपनगरम्}} (\emph{upanagaram}) is built from \emph{upa}, \textquotedblleft{}near\textquotedblright{}, and \emph{nagara}, \textquotedblleft{}a city\textquotedblright{}; MW gives the compound as a suburb (MW).
\end{cousinweb}

\begin{grammarlens}[title={What you've built}]
One more word for getting around.
\end{grammarlens}

\subsection*{Your turn}
\begin{itemize}
\raggedright
  \item \emph{Say it:} \emph{upanagaram}
  \item \emph{Say it:} \emph{upanagaram}, once more
  \item \emph{Say it:} \emph{mama mitram upanagare vasati.} — My friend lives in a suburb.
  \item \emph{Recall:} say the Sanskrit for refreshes, then the Sanskrit for a hamlet, then say \emph{upanagaram} again
\end{itemize}

\begin{tcolorbox}[breakable,colback=teal!4,colframe=teal!35!black,title={Before you move on}]
What does \sk{उपनगरम्} mean? (\textquotedblleft{}A suburb\textquotedblright{}.) Say it once more.
\end{tcolorbox}
\section[\texorpdfstring{\sk{चत्वरम्} (catvaram)}{चत्वरम् (catvaram)}]{\sk{चत्वरम्} (catvaram) — a town square}

\label{lesson:SA-C421-catvaram}



\begin{quote}
Before the new one: say the Sanskrit for a hamlet, then the Sanskrit for a suburb.
\end{quote}

\subsection*{You'll want to know: \sk{चत्वरम्}}
\textbf{\sk{चत्वरम्}} — \emph{catvaram} — \textquotedblleft{}a town square\textquotedblright{}.

\begin{quote}
\sk{सायं} \sk{चत्वरे} \sk{बालाः} \sk{क्रीडन्ति।}
\end{quote}

\begin{quote}
\emph{sāyaṃ catvare bālāḥ krīḍanti.}
\end{quote}

\begin{quote}
In the evening children play in the square.
\end{quote}

A neuter noun.

\begin{cousinweb}
\textbf{\sk{चत्वरम्}} (\emph{catvaram}) is the word MW gives for a square place, a place where many ways meet; MW gives no derivation, so its origin is uncertain (MW).
\end{cousinweb}

\begin{grammarlens}[title={What you've built}]
One more word for getting around.
\end{grammarlens}

\subsection*{Your turn}
\begin{itemize}
\raggedright
  \item \emph{Say it:} \emph{catvaram}
  \item \emph{Say it:} \emph{catvaram}, once more
  \item \emph{Say it:} \emph{sāyaṃ catvare bālāḥ krīḍanti.} — In the evening children play in the square.
  \item \emph{Recall:} say the Sanskrit for a hamlet, then the Sanskrit for a suburb, then say \emph{catvaram} again
\end{itemize}

\begin{tcolorbox}[breakable,colback=teal!4,colframe=teal!35!black,title={Before you move on}]
What does \sk{चत्वरम्} mean? (\textquotedblleft{}A town square\textquotedblright{}.) Say it once more.
\end{tcolorbox}
\section[\texorpdfstring{\sk{हर्म्यम्} (harmyam)}{हर्म्यम् (harmyam)}]{\sk{हर्म्यम्} (harmyam) — a mansion}

\label{lesson:SA-C421-harmyam}



\begin{quote}
Before the new one: say the Sanskrit for a suburb, then the Sanskrit for a town square.
\end{quote}

\subsection*{You'll want to know: \sk{हर्म्यम्}}
\textbf{\sk{हर्म्यम्}} — \emph{harmyam} — \textquotedblleft{}a mansion\textquotedblright{}.

\begin{quote}
\sk{नगरे} \sk{धनिकस्य} \sk{विशालं} \sk{हर्म्यम्} \sk{अस्ति।}
\end{quote}

\begin{quote}
\emph{nagare dhanikasya viśālaṃ harmyam asti.}
\end{quote}

\begin{quote}
In the city there is a rich man's huge mansion.
\end{quote}

A neuter noun.

\begin{cousinweb}
\textbf{\sk{हर्म्यम्}} (\emph{harmyam}) is the word MW gives for a large house, a palace, a mansion; MW doubts the old derivation from √hṛ and connects it rather with \emph{gharma}, \textquotedblleft{}heat\textquotedblright{}, so its origin is uncertain (MW).
\end{cousinweb}

\begin{grammarlens}[title={What you've built}]
One more word for getting around.
\end{grammarlens}

\subsection*{Your turn}
\begin{itemize}
\raggedright
  \item \emph{Say it:} \emph{harmyam}
  \item \emph{Say it:} \emph{harmyam}, once more
  \item \emph{Say it:} \emph{nagare dhanikasya viśālaṃ harmyam asti.} — In the city there is a rich man's huge mansion.
  \item \emph{Recall:} say the Sanskrit for a suburb, then the Sanskrit for a town square, then say \emph{harmyam} again
\end{itemize}

\begin{tcolorbox}[breakable,colback=teal!4,colframe=teal!35!black,title={Before you move on}]
What does \sk{हर्म्यम्} mean? (\textquotedblleft{}A mansion\textquotedblright{}.) Say it once more.
\end{tcolorbox}
\section[\texorpdfstring{\sk{सम्भाषते} (sambhāṣate)}{सम्भाषते (sambhāṣate)}]{\sk{सम्भाषते} (sambhāṣate) — converses; chats}

\label{lesson:SA-C421-sambhasate}



\begin{quote}
Before the new one: say the Sanskrit for a town square, then the Sanskrit for a mansion.
\end{quote}

\subsection*{You'll want to know: \sk{सम्भाषते}}
\textbf{\sk{सम्भाषते}} — \emph{sambhāṣate} — \textquotedblleft{}converses; chats\textquotedblright{}.

\begin{quote}
\sk{यात्री} \sk{चत्वरे} \sk{जनैः} \sk{सह} \sk{सम्भाषते।}
\end{quote}

\begin{quote}
\emph{yātrī catvare janaiḥ saha sambhāṣate.}
\end{quote}

\begin{quote}
The traveller chats with people in the square.
\end{quote}

Past (he …): \sk{सः} \sk{जनैः} \sk{सह} \sk{सम्भाषितवान्}, \emph{saḥ janaiḥ saha sambhāṣitavān}.

\begin{cousinweb}
\textbf{\sk{सम्भाषते}} (\emph{sambhāṣate}) is built from \emph{sam}, \textquotedblleft{}together\textquotedblright{}, and the root √bhāṣ, \textquotedblleft{}to speak\textquotedblright{}, which also gives \emph{bhāṣate}, \textquotedblleft{}speaks\textquotedblright{}; MW gives \emph{sam-bhāṣ} as to speak together, to converse (MW).
\end{cousinweb}

\begin{grammarlens}[title={What you've built}]
One more word for getting around.
\end{grammarlens}

\subsection*{Your turn}
\begin{itemize}
\raggedright
  \item \emph{Say it:} \emph{sambhāṣate}
  \item \emph{Say it:} \emph{sambhāṣate}, once more
  \item \emph{Say it:} \emph{yātrī catvare janaiḥ saha sambhāṣate.} — The traveller chats with people in the square.
  \item \emph{Recall:} say the Sanskrit for a town square, then the Sanskrit for a mansion, then say \emph{sambhāṣate} again
\end{itemize}

\begin{tcolorbox}[breakable,colback=teal!4,colframe=teal!35!black,title={Before you move on}]
What does \sk{सम्भाषते} mean? (\textquotedblleft{}Converses; chats\textquotedblright{}.) Say it once more.
\end{tcolorbox}
